\section{Amalgams and primary torsion}\label{amalgams-and-primary-torsion}

Fix injective homomorphisms \(i_A:C\to A\) and \(i_B:C\to B\). Let
\(A*B\) be the free product and let \[
P=(A*B)/\langle\!\langle i_A(c)i_B(c)^{-1}:c\in C\rangle\!\rangle.
\tag{A1}
\] We use the amalgam normal-form theorem in the following precise form:
the canonical maps of \(A\) and \(B\) into \(P\) are injective; an
alternating reduced word with factors outside the common copy of \(C\)
has nonzero reduced length and is nonidentity; every element is
conjugate either into a factor or to a cyclically reduced alternating
word of length at least two, whose positive powers remain reduced with
proportionally increasing length. This is a retained classical theorem.

\textbf{Theorem UA.Amalgam (universal property and torsion).} The group
\(P\) in (A1), with its canonical maps \(j_A,j_B\), is a pushout of
\(A\leftarrow C\to B\) in the category of groups: for every group \(D\)
and homomorphisms \(f_A:A\to D\), \(f_B:B\to D\) satisfying
\(f_Ai_A=f_Bi_B\), there is a unique homomorphism \(f:P\to D\) with
\(fj_A=f_A\) and \(fj_B=f_B\).

Let \(p\) be a prime. Say that a group has \(p\)-primary finite torsion
when every finite-order element has \(p\)-power order. This property
passes to subgroups, to extensions with finite \(p\)-group quotient, and
to amalgamated free products of groups having the property. If
\(i_A(C)\ne A\) and \(i_B(C)\ne B\), then \(P\) is infinite. If \(A\)
and \(B\) are finitely generated, then so is \(P\).

\emph{Proof.} The free-product universal property gives a unique
homomorphism \(A*B\to D\) extending \(f_A,f_B\). Agreement on \(C\) puts
every defining relator in its kernel, so it factors through \(P\). The
images of \(A\) and \(B\) generate \(P\); hence this factorization is
unique. The equality \(j_Ai_A=j_Bi_B\) follows from the defining
relations, completing the pushout property.

The torsion assertion for subgroups is immediate. For an extension
\(1\to K\to E\to R\to1\), with \(R\) a finite \(p\)-group and \(K\)
having \(p\)-primary finite torsion, take any finite cyclic subgroup
\(\langle x\rangle\) of \(E\). Its intersection with \(K\) is a cyclic
\(p\)-group, and its image in \(R\) is a \(p\)-group. The order formula
for this restricted exact sequence makes \(|x|\) a power of \(p\).

For an amalgam, the normal-form theorem places a torsion element, up to
conjugacy, in a factor: the alternative cyclically reduced word has
powers of unbounded reduced length and cannot have finite order. The
factor hypothesis therefore proves the torsion assertion. If the common
subgroup is proper in both factors, choose \(a\in A\setminus i_A(C)\)
and \(b\in B\setminus i_B(C)\). The words \((ab)^n\) are reduced of
length \(2n\), so are all distinct; this proves infinitude. Finally, the
images of finite generating sets of \(A\) and \(B\) generate their free
product and its quotient \(P\), proving finite generation. \(\square\)

The simultaneous fiber product of UA.SK has its displayed compatibility
equations and its proved generation theorem. It is not the pushout in
UA.Amalgam. A graph of proof dependencies is neither of these group
constructions.
